% Apéndice: rúbricas de autoevaluación
\section{Rúbricas de autoevaluación}\label{ap:rubricas}

\textbf{Material docente.} Al final de la clase 1 el profesor armó con los alumnos ``una
especie de rúbrica para autoevaluarnos''. Se reproduce textual (sólo se corrigió
ortografía). Para los demás bloques del apunte se propone una rúbrica análoga, redactada
por la ayudantía a partir de los resúmenes de clase; no es del profesor. En todas, cada
casillero nombra la propiedad algebraica \emph{y} el hecho geométrico que la explica.

\subsection{Rotaciones (rúbrica del profesor, clase 1)}

{\small
\begin{tabularx}{\textwidth}{YYY}
\toprule
\textbf{Elemental} & \textbf{Medio} & \textbf{Avanzado}\\
\midrule
Las matrices y las transformaciones lineales tienen una correspondencia 1 a 1 &
Combinación de rotaciones son rotaciones &
Existe la derivada de una transformación, y genera un conjunto de vectores tangentes a
trayectorias\\
\addlinespace
Las rotaciones son matrices ortogonales (\(A^{T}A=Id\)) con determinante mayor a 0 &
Hay transformaciones que se definen completamente por qué propiedades de los objetos
transformados mantienen &
Cuando hay estructuras comunes entre dos conjuntos de elementos y cómo son afectados por
distintas transformaciones (una cada conjunto) hay una equivalencia que se puede definir\\
\addlinespace
Multiplicar por \(i\) es rotar \(\pi/2\) &
La exponencial (por cómo se deriva) juega un rol especial cuando hablamos de
transformaciones &
Un recurso útil para analizar transformaciones es lo que definimos coloquialmente como
``ir a lo chico para volver a lo grande''\\
\addlinespace
Derivar funciones con números complejos sigue las mismas reglas que los reales &
Hay una equivalencia entre operaciones en \(\mathbb{C}\) y en \(\R^2\) (que se describe
por lo que ``le hacen'' a los puntos que intervienen) & \\
\addlinespace
La falta de conmutatividad de transformaciones lineales depende de la magnitud &
Las transformaciones se identifican con las trayectorias que generan (por aplicación
repetida) & \\
\addlinespace
El producto interno puede servir como base para definir tipos de transformaciones &
Descomponer una matriz en ``la identidad + algo'' nos da una pista de su comportamiento
diferencial si el algo es chiquito & \\
\bottomrule
\end{tabularx}
}

\subsection{Afinidades y coordenadas homogéneas (propuesta de la ayudantía)}

{\small
\begin{tabularx}{\textwidth}{YYY}
\toprule
\textbf{Elemental} & \textbf{Medio} & \textbf{Avanzado}\\
\midrule
Una afinidad conserva rectas las rectas y paralelas las paralelas; en coordenadas es
\(F(x)=Ax+t\). Hecho geométrico: un cuadrado va a un paralelogramo, nunca a un trapecio &
Dos reflexiones en rectas secantes son una rotación (conservan distancias y dejan un punto
fijo); con reflexiones y escalas axiales se construye cualquier afinidad &
Existe una única afinidad que lleva un triángulo dado a otro triángulo dado; saber
construirla gráficamente (Pamfilos) y algebraicamente (tres puntos, seis ecuaciones)\\
\addlinespace
En coordenadas homogéneas la traslación entra en la matriz \(3\times3\) con última fila
\((0,0,1)\). Hecho geométrico: el plano afín es el plano \(w=1\) y cada punto es un rayo &
Componer afinidades es multiplicar matrices \(3\times3\); el orden importa salvo casos
especiales (mismo centro, dos traslaciones). Hecho geométrico: rotar y luego trasladar no
termina donde trasladar y luego rotar &
Decidir cuándo dos afinidades conmutan y por qué (comparten centro, ejes o dirección), y
qué figura describe la composición repetida\\
\addlinespace
Un punto fijo cumple \(Ax+t=x\). Hecho geométrico: es el lugar donde la figura se
``ancla'' &
Rotar y trasladar sucesivamente un punto dibuja círculos alrededor del punto fijo de la
composición; hallar ese centro &
Los puntos fijos del plano son los autovectores de la matriz \(3\times3\); un autovector
con tercera coordenada nula es una dirección conservada (punto ideal fijo)\\
\addlinespace
Aplicar una afinidad en planilla: matriz por vector columna, punto por punto. Hecho
geométrico: las columnas son las imágenes de los ejes &
Combinando rotaciones y una matriz diagonal se le ponen coordenadas a cualquier afinidad.
Hecho geométrico: girar, estirar en dos direcciones perpendiculares, volver a girar &
Ver esa descomposición rotación--diagonal--rotación como dos afinidades axiales sucesivas
y reconocer en ella el adelanto de la SVD\\
\bottomrule
\end{tabularx}
}

\subsection{Transformaciones proyectivas (propuesta de la ayudantía)}

{\small
\begin{tabularx}{\textwidth}{YYY}
\toprule
\textbf{Elemental} & \textbf{Medio} & \textbf{Avanzado}\\
\midrule
Una proyectiva conserva rectas, nada más; en homogéneas es cualquier matriz \(3\times3\)
invertible. Hecho geométrico: la foto de un lote rectangular es un trapecio &
Cuatro puntos, sin tres alineados, determinan la homografía; la afinidad es el caso con
última fila \((0,0,1)\). Hecho geométrico: con tres puntos falta decidir ``cuánta
perspectiva hay'' &
La recta \(W=0\) se va al infinito (el horizonte de la foto); puntos ideales que se vuelven
propios y viceversa; qué pasa con una figura que cruza esa recta\\
\addlinespace
Los puntos son rayos: \((x,y,1)\) y \(\lambda(x,y,1)\) son el mismo punto; normalizar
dividiendo por \(W\) recupera el plano &
Los puntos fijos son los autovectores de \(H\) y siempre hay al menos uno real. Hecho
geométrico: toda homografía deja quieto algún punto, propio o ideal &
Las rectas fijas son formas lineales fijas: dualidad punto--recta (\(\ell\cdot p=0\)) y su
lectura como autovectores de \(H^{T}\)\\
\addlinespace
Las paralelas dejan de serlo: se cortan en un punto de fuga. Hecho geométrico: el punto
ideal común se hizo propio &
Jerarquía euclídeas \(\subset\) similaridades \(\subset\) afines \(\subset\) proyectivas, y
qué invariante se pierde en cada paso (distancia, ángulo, paralelismo, razón simple) &
Razón doble como el invariante que sobrevive; reconocer la homología (centro sobre el eje)
y sus autovalores repetidos\\
\addlinespace
Reconocer una proyectiva en una imagen y ubicar el punto de fuga &
Calcular la homografía de un cuadrilátero con el patrón de cuatro puntos y verificar
\(Hp_i=s_iq_i\) &
Analizar el caso en que la imagen de un polígono ``se parte'' porque un lado cruza la recta
que va al infinito (ejercicio del pentágono)\\
\bottomrule
\end{tabularx}
}

\subsection{Productos, determinante y autovectores (propuesta de la ayudantía)}

{\small
\begin{tabularx}{\textwidth}{YYY}
\toprule
\textbf{Elemental} & \textbf{Medio} & \textbf{Avanzado}\\
\midrule
El producto interno decide incidencia y ortogonalidad; el vectorial une dos puntos y corta
dos rectas. Hecho geométrico: \(\ell\cdot p=0\) es ``el punto está en la recta'' &
El determinante es la única operación lineal en cada argumento y antisimétrica; vale cero
si hay dependencia lineal. Hecho geométrico: es el área o volumen orientado &
El producto cuña generaliza el determinante a cualquier dimensión y ``arma formas
lineales'' que detectan colinealidad; ortogonal es lo que maximiza el volumen\\
\addlinespace
Una forma lineal es una regla graduada: familia de rectas paralelas equiespaciadas; la
matriz es una pila de reglas &
Las formas lineales forman un espacio vectorial (el dual) y se transforman del lado
contrario que los vectores. Hecho geométrico: las rectas de nivel se arrastran con la figura &
El gradiente es una forma lineal cuya recta de nivel cero es la tangente; de ahí que sea
perpendicular a la curva de nivel\\
\addlinespace
Un autovector cumple \(Av=\lambda v\). Hecho geométrico: es una dirección que cae sobre sí
misma, sólo se estira o se invierte &
\(Av=\lambda v\Leftrightarrow(A-\lambda I)v=0\Leftrightarrow\det(A-\lambda I)=0\); el
autovector sale del producto cruz de dos filas. Hecho geométrico: es el punto común de las
rectas de nivel cero de las filas &
Autovalor doble (todo un plano de autovectores), núcleo previo, autovalores complejos (una
rotación escondida): qué figura corresponde a cada caso\\
\addlinespace
Una transformación lineal tiene una única matriz y la composición es el producto. Hecho
geométrico: las columnas dicen adónde van los ejes &
El núcleo es lo que va al cero; el cero siempre va al cero. Hecho geométrico: una
transformación con núcleo aplasta una dirección entera a un punto &
Los cuatro subespacios de una transformación (núcleo, imagen, co-imagen, espacio
inaccesible de la salida) y su dibujo para una matriz \(2\times2\) de rango 1\\
\addlinespace
El eje de una rotación 3D es el autovector de autovalor \(1\) &
Los puntos fijos de una afinidad o de una proyectiva son subespacios invariantes en
\(\R^3\) &
Construir una transformación con un vector dado en el núcleo usando el producto cuña, y
saber qué le hace al resto del espacio\\
\bottomrule
\end{tabularx}
}

\subsection{Cambio de base y SVD (propuesta de la ayudantía)}

{\small
\begin{tabularx}{\textwidth}{YYY}
\toprule
\textbf{Elemental} & \textbf{Medio} & \textbf{Avanzado}\\
\midrule
Un cambio de base expresa el mismo vector como combinación de otros vectores; la matriz
tiene en sus columnas una base escrita en la otra. Hecho geométrico: la misma flecha medida
con otra grilla &
Cambiar la base de la matriz de una transformación es componer con los dos cambios de
base: \(M_{av}=C_{av\leftarrow e}\,M\,C_{e\leftarrow av}\). Hecho geométrico: traducir,
transformar, volver a traducir &
No hay base privilegiada a priori; lo único especial de la canónica es que en ella el
producto interno se calcula como suma de productos. Hecho geométrico: sin distancias, todas
las grillas valen lo mismo\\
\addlinespace
Si hay tantos autovectores independientes como dimensiones, en esa base la matriz es
diagonal. Hecho geométrico: la transformación es un estiramiento en cada dirección propia &
La base de autovectores permite potencias enteras y reales de la matriz. Hecho geométrico:
ver la transformación como un proceso continuo y dibujar la trayectoria de un punto &
Autovalores complejos: por qué la diagonalización real puede fallar. Hecho geométrico: hay
un giro escondido y ninguna dirección queda quieta\\
\addlinespace
Existe una base ortogonal de la entrada que va a una base ortogonal de la salida. Hecho
geométrico: la circunferencia unidad se transforma en una elipse y sus ejes son esa base &
En la SVD las bases de entrada y salida son distintas, \(A=U\Sigma V^{T}\), y la diagonal
nunca es compleja. Hecho geométrico: girar, estirar en direcciones perpendiculares, girar &
Matrices simétricas: si un autovector coincide con un vector singular, coinciden todos (en
\(\R^2\)). Preguntarse si pasa lo mismo en \(\R^3\)\\
\addlinespace
Invertir una matriz diagonal es tomar recíprocos. Hecho geométrico: deshacer cada
estiramiento &
La inversa es siempre ``respecto de una operación'': el elemento que compuesto con otro da
el neutro (suma, producto, composición, producto de matrices) &
Pseudoinversa de Moore--Penrose: resolver un sistema sobredeterminado eligiendo el punto de
la imagen más cercano al término independiente. Hecho geométrico: proyectar sobre la imagen\\
\bottomrule
\end{tabularx}
}

